\section*{Bảng phân công nhiệm vụ.}
\addcontentsline{toc}{section}{Bảng phân công nhiệm vụ.}

\begin{table}[htbp]
    \centering
    \renewcommand{\arraystretch}{1.4} % Tăng độ thoáng theo chiều dọc cho các hàng
    % Tổng chiều rộng bảng = \linewidth (nằm gọn trong lề trang giấy)
    \begin{tabularx}{\linewidth}{|C{2.2cm}|C{4.5cm}|X|}
        \hline
        \textbf{MSSV} & \multicolumn{1}{c|}{\textbf{Họ và tên}} & \multicolumn{1}{c|}{\textbf{Nhiệm vụ}} \\
        \hline
        
        % Hàng 01.
        25120269 & 
        Nguyễn Hoàng Quân \newline (Nhóm trưởng) & 
        \begin{tabitemize}
            \item Phân công nhiệm vụ cho các thành viên.
            \item Kiểm tra lại nội dung lý thuyết.
            \item Tổng hợp nội dung báo cáo.
            \item \textbf{Câu 03:} Tìm hiểu về Composite Design Pattern trong OOP (phần code).
        \end{tabitemize} \\
        \hline

        % Hàng 02.
        25120267 & 
        Văn Thành Phát & 
        \begin{tabitemize}
            \item \textbf{Câu 02:} Tìm hiểu về function template và class template trong C++.
            \item Chỉnh sửa, format Latex và xuất ra file PDF.
        \end{tabitemize} \\
        \hline
        
        % Hàng 03.
        25120264 & 
        Võ Ngọc Linh & 
        \begin{tabitemize}
            \item \textbf{Câu 01:} Tìm hiểu về đa kế thừa trong C++ (vấn đề hình thoi và cách giải quyết).
        \end{tabitemize} \\
        \hline
        
        % Hàng 04.
        25120271 & 
        Lê Nguyễn Như Quỳnh & 
        \begin{tabitemize}
            \item \textbf{Câu 03:} Tìm hiểu về Composite Design Pattern trong OOP (phần lý thuyết).
        \end{tabitemize} \\
        \hline
        
        % Hàng 05.
        25120265 & 
        Trương Huỳnh Như & 
        \begin{tabitemize}
            \item \textbf{Câu 01:} Tìm hiểu về đa kế thừa trong C++ (khái niệm về đa kế thừa, vấn đề trùng tên thuộc tính, vấn đề trùng tên hàm và cách giải quyết).
        \end{tabitemize} \\
        \hline
    \end{tabularx}
\end{table}